\chapter{Requisitos del sistema}
\label{chap:requisitos}

\lettrine{X}{xxx},\dots

\section{Introducción}

{ \color{red} Aspectos a considerar
    \begin{itemize}
      \item Transcripción textual/completa de los requisitos, punto de partida para identificar actores y a continuación casos de uso (siguientes secciones). 
      \item Pueden diferenciarse entre requisitos funcionales y no funcionales
    \end{itemize}
}

\section{Actores}

\section{Casos de uso}

{ \color{red} Aspectos a considerar
    \begin{itemize}
      \item Agrupar casos de uso (servirá de base para agruparlos en servicios posteriormente en la sección de diseño).
      \item Definición concreta /detallada de cada caso de uso (precondiciones ...)
    \end{itemize}
}

\section{Modelo de casos de uso}

{ \color{red} Aspectos a considerar
    \begin{itemize}
      \item Diagrama de casos de uso incluyendo actores y relaciones entre casos
    \end{itemize}
}

\section{Prototipado de interfaz}

{ \color{red} Es interesante hacer el esfuerzo en este momento de prototipar las principales pantallas de la interfaz (y añadirlas en la memoria), como parte del análisis de requisitos, así como de aspectos de usabilidad de las mismas. }